% WebTKG-RAG: Structure-Aware Retrieval-Augmented Generation
% over Evolving Web Documents via DOM-Native Temporal Knowledge Graphs
%
% Target: WWW 2026 / ACL 2026 / arXiv preprint
% Format: ACL 2024 style (also compatible with arXiv)

\documentclass[11pt]{article}
\usepackage[utf8]{inputenc}
\usepackage{amsmath,amssymb,amsfonts}
\usepackage{graphicx}
\usepackage{booktabs}
\usepackage{multirow}
\usepackage{hyperref}
\usepackage{algorithm}
\usepackage{algorithmic}
\usepackage{xcolor}
\usepackage{subcaption}
\usepackage[margin=1in]{geometry}

\newcommand{\ours}{\textsc{WebTKG-RAG}}
\newcommand{\todo}[1]{\textcolor{red}{[TODO: #1]}}

\title{\ours{}: Structure-Aware Retrieval-Augmented Generation \\
over Evolving Web Documents via \\
DOM-Native Temporal Knowledge Graphs}

\author{
  [Author Names] \\
  [Affiliations] \\
  \texttt{[emails]}
}

\date{}

\begin{document}
\maketitle

% ============================================================
\begin{abstract}
Retrieval-Augmented Generation (RAG) systems increasingly rely on web content as external knowledge, yet current approaches discard the rich structural and temporal information inherent in HTML documents. Existing methods either convert HTML to plain text---losing structural semantics---or treat HTML as a static snapshot, ignoring how web content evolves over time. We propose \ours{}, a novel framework that treats the Document Object Model (DOM) tree as a first-class, temporally-evolving, visually-grounded knowledge graph for RAG. Our approach introduces three key innovations: (1)~\textbf{tri-modal DOM node embeddings} that fuse textual content, tree-structural position, and visual rendering layout; (2)~a \textbf{temporal DOM knowledge graph} constructed via structural tree differencing across time snapshots, enabling time-aware question answering; and (3)~\textbf{confidence-guided tree traversal} for retrieval, which navigates the DOM hierarchy top-down with learned relevance scores, achieving sub-linear retrieval complexity. We further propose \textbf{cross-site DOM fingerprinting} via contrastive learning, enabling zero-shot structured extraction on unseen websites. Experiments on six QA benchmarks, two web extraction datasets, and a new temporal web QA benchmark demonstrate that \ours{} outperforms state-of-the-art methods including HtmlRAG, achieving X\% improvement on static QA and Y\% on temporal QA tasks, while reducing retrieval token count by Z\%.
\end{abstract}

% ============================================================
\section{Introduction}
\label{sec:intro}

The Web is the largest and most dynamic knowledge source available to Retrieval-Augmented Generation (RAG) systems~\cite{htmlrag2025,rag_survey2025}. Commercial systems such as Perplexity and SearchGPT retrieve web pages, extract their content, and feed it to large language models (LLMs) to generate grounded answers. However, a fundamental disconnect exists between how web content is \emph{authored} and how it is \emph{consumed} by RAG systems.

Web pages are authored as HTML---a hierarchical markup language that encodes rich structural semantics through the Document Object Model (DOM) tree. A product page, for instance, organizes information into semantically meaningful regions: product titles in \texttt{<h1>} tags, prices in \texttt{<span class="price">} elements, availability indicators in specific \texttt{<div>} containers, and images in \texttt{<img>} tags. This structure is not arbitrary; it reflects the \emph{semantic ontology} of the underlying data.

Yet current RAG pipelines systematically destroy this structure. The dominant paradigm---``Retrieve HTML $\rightarrow$ Convert to Plain Text $\rightarrow$ Chunk $\rightarrow$ Embed $\rightarrow$ Retrieve $\rightarrow$ Generate''---strips away the DOM hierarchy, discards visual layout information, and treats the resulting flat text as if it were a prose document. Recent work by Tan et al.~\cite{htmlrag2025} demonstrated that preserving HTML format improves RAG performance, but their approach still treats HTML as enhanced text rather than as a native knowledge structure.

We identify three critical limitations of existing approaches:

\begin{enumerate}
    \item \textbf{Structural Blindness.} Converting HTML to text or even preserving tags loses the \emph{spatial} and \emph{hierarchical} relationships between elements. A price displayed prominently at the center of a product page carries different semantic weight than the same number buried in a footer, but text-based embeddings cannot distinguish them.

    \item \textbf{Temporal Amnesia.} Web content is inherently dynamic---prices change, products go in and out of stock, sales begin and end. Current RAG systems operate on static snapshots, unable to answer temporal queries like ``When did this product's price drop?'' or ``Has this item been on sale recently?''

    \item \textbf{Site-Specific Fragility.} Extraction and retrieval methods trained on one website's DOM structure fail when applied to another, even when both encode the same semantic information (e.g., product price) in structurally different HTML.
\end{enumerate}

To address these limitations, we propose \ours{} (\textbf{Web} \textbf{T}emporal \textbf{K}nowledge \textbf{G}raph \textbf{RAG}), a framework that reconceptualizes the relationship between HTML and knowledge. Our key insight is that \emph{the DOM tree itself is a knowledge graph}---nodes represent semantic entities, edges encode structural relationships, and the temporal evolution of this graph captures how knowledge changes over time.

\ours{} makes the following contributions:

\begin{itemize}
    \item We propose \textbf{tri-modal DOM node embeddings} that jointly encode textual content, tree-structural position (depth, sibling index, tag type, XPath), and visual rendering properties (bounding box coordinates from headless browser rendering), creating a unified representation that captures what information a node contains, where it sits in the document hierarchy, and how it appears visually.

    \item We introduce a \textbf{temporal DOM knowledge graph} constructed by computing structural tree diffs (via the Zhang-Shasha algorithm) across time-ordered DOM snapshots. Changes are converted into timestamped knowledge triples, enabling the RAG system to answer temporal queries with full provenance.

    \item We design a \textbf{confidence-guided tree traversal} algorithm for retrieval that navigates the DOM hierarchy top-down, pruning irrelevant subtrees early and drilling into promising ones, achieving sub-linear retrieval complexity compared to flat top-$k$ approaches.

    \item We propose \textbf{cross-site DOM fingerprinting} via contrastive learning on DOM subtree structures, enabling zero-shot structured extraction on previously unseen websites.

    \item We create and release a new \textbf{Temporal Web QA} benchmark consisting of product pages crawled over multiple time points with annotated temporal questions.
\end{itemize}

% ============================================================
\section{Related Work}
\label{sec:related}

\subsection{HTML-Aware RAG Systems}

Traditional RAG systems convert retrieved HTML documents to plain text before processing~\cite{langchain2022,llamaindex2022}. HtmlRAG~\cite{htmlrag2025} was the first to argue that HTML format preserves richer semantics than plain text for RAG, proposing a block-tree-based pruning strategy that reduces HTML length while retaining key information. However, HtmlRAG treats HTML purely as text with structural tags---it does not leverage visual rendering information, does not model temporal evolution, and does not support cross-site transfer. Our work extends this direction by treating the DOM tree as a first-class knowledge graph with tri-modal node representations.

\subsection{DOM Tree Representation Learning}

Several works have explored learning representations from DOM trees. DOM-LM~\cite{domlm2022} pre-trains language models on HTML documents with tree position awareness. MarkupLM~\cite{markuplm2022} jointly pre-trains text and markup language using XPath embeddings. WebFormer~\cite{webformer2022} designs rich attention mechanisms between HTML tokens and DOM nodes. More recently, the Klarna Product Page Dataset~\cite{klarna2024} benchmarked Graph Neural Networks (GCN, GAT, GraphSAGE) on DOM trees for web element nomination. Our work differs by (1) adding visual rendering as a third modality, (2) applying DOM representations to RAG retrieval rather than classification, and (3) modeling temporal evolution.

\subsection{Web Information Extraction}

AXE~\cite{axe2025} treats the HTML DOM as a tree to prune, achieving 88.1\% F1 on the SWDE dataset with a tiny 0.6B LLM. SCRIBES~\cite{scribes2025} uses reinforcement learning to generate reusable extraction scripts for structurally similar pages. Pinterest's cross-domain web IE system~\cite{pinterest2025} combines structural, visual, and text modalities per HTML node---the closest work to our tri-modal approach, but applied to extraction rather than RAG retrieval. Prune4Web~\cite{prune4web2025} generates Python scoring programs for DOM element filtering. Our work unifies extraction and retrieval in a single framework through the DOM-as-knowledge-graph paradigm.

\subsection{Visual Document Retrieval}

ColPali~\cite{colpali2024} introduced late-interaction visual retrieval using vision-language model patch embeddings, treating documents as images. Snappy~\cite{snappy2024} extends this with patch-to-region relevance propagation. SCAN~\cite{scan2025} performs semantic document layout analysis for RAG. These approaches target PDF/scanned documents and operate on rendered images. Our approach is complementary: we leverage the DOM tree structure \emph{and} visual rendering, obtaining spatial information directly from the browser's layout engine rather than from image patches.

\subsection{Temporal Knowledge Graphs and RAG}

TG-RAG~\cite{tgrag2025} models external corpora as bi-level temporal graphs with timestamped relations. STAR-RAG~\cite{starrag2025} builds time-aligned rule graphs. StreamingRAG~\cite{streamingrag2025} constructs evolving knowledge graphs in real-time from video streams. EvoKG~\cite{evokg2025} incrementally updates knowledge graphs from unstructured documents. ATOM~\cite{atom2025} splits documents into atomic facts for temporal KG construction. All these works operate on text corpora or multimedia streams---none apply temporal knowledge graph construction to HTML DOM trees, which is our key contribution.

\subsection{DOM Change Detection}

HDNA~\cite{hdna2023} proposes graph-based change detection in HTML pages using weighted DOM tree comparison. The Zhang-Shasha algorithm~\cite{zhangshasha1989} provides the foundational $O(n^2)$ tree edit distance computation. VDiff~\cite{vdiff2004} applies tree edit distance specifically to HTML differencing. We build upon these algorithms but go beyond detection: we convert detected changes into temporal knowledge triples that power a RAG system.

% ============================================================
\section{Methodology}
\label{sec:method}

\subsection{Overview}

Figure~\ref{fig:overview} presents the \ours{} pipeline. Given a collection of HTML documents (potentially crawled at multiple time points), our system: (1) parses each document into a DOM tree and constructs a DOM Knowledge Graph with tri-modal node embeddings (\S\ref{sec:trimodal}); (2) computes temporal diffs between time-ordered snapshots to build a Temporal DOM Knowledge Graph (\S\ref{sec:temporal}); (3) indexes the knowledge graph for retrieval using confidence-guided tree traversal (\S\ref{sec:retrieval}); and (4) feeds retrieved HTML context to an LLM for answer generation.

\subsection{DOM Tree as Knowledge Graph}
\label{sec:domkg}

We formalize the DOM tree as a knowledge graph $\mathcal{G} = (\mathcal{V}, \mathcal{E}, \mathcal{A})$ where:
\begin{itemize}
    \item $\mathcal{V}$ is the set of DOM nodes (elements, text nodes)
    \item $\mathcal{E} \subseteq \mathcal{V} \times \mathcal{V}$ encodes parent-child and sibling relationships
    \item $\mathcal{A}: \mathcal{V} \rightarrow \mathbb{R}^d$ maps each node to a tri-modal embedding
\end{itemize}

Unlike HtmlRAG's block tree, which merges DOM nodes into coarse blocks for pruning, our knowledge graph preserves the full DOM structure while augmenting each node with rich multi-modal features.

\subsection{Tri-Modal DOM Node Embeddings}
\label{sec:trimodal}

For each DOM node $v \in \mathcal{V}$, we compute three modality-specific representations and fuse them into a unified embedding.

\paragraph{Textual Embedding $\mathbf{e}_v^{\text{text}}$.}
We extract the text content of node $v$ (including direct text and alt attributes) and encode it using a pre-trained text encoder (e.g., BGE-Large):
\begin{equation}
    \mathbf{e}_v^{\text{text}} = \text{TextEncoder}(\text{content}(v))
\end{equation}

\paragraph{Structural Embedding $\mathbf{e}_v^{\text{struct}}$.}
We encode the node's position in the DOM tree using:
\begin{itemize}
    \item \textbf{Tag type}: One-hot encoding of the HTML tag (e.g., \texttt{div}, \texttt{span}, \texttt{h1}, \texttt{p})
    \item \textbf{Tree depth}: Normalized depth in the DOM tree
    \item \textbf{Sibling index}: Position among siblings
    \item \textbf{XPath encoding}: Following MarkupLM~\cite{markuplm2022}, we encode the XPath from root to node as a sequence of tag embeddings
\end{itemize}
\begin{equation}
    \mathbf{e}_v^{\text{struct}} = \text{StructEncoder}(\text{tag}(v), \text{depth}(v), \text{sibling\_idx}(v), \text{xpath}(v))
\end{equation}

\paragraph{Visual Embedding $\mathbf{e}_v^{\text{visual}}$.}
We render the HTML document in a headless browser (Playwright) and extract the bounding box $(x, y, w, h)$ for each DOM node via \texttt{getBoundingClientRect()}. We also compute:
\begin{itemize}
    \item \textbf{Relative position}: $(x/W, y/H, w/W, h/H)$ normalized by viewport dimensions
    \item \textbf{Visual area}: $w \times h$ as a proxy for visual prominence
    \item \textbf{Font size}: Computed CSS font-size property
    \item \textbf{Color contrast}: Foreground-background color difference
\end{itemize}
\begin{equation}
    \mathbf{e}_v^{\text{visual}} = \text{VisualEncoder}(\text{bbox}(v), \text{area}(v), \text{font}(v), \text{contrast}(v))
\end{equation}

\paragraph{Fusion.}
We fuse the three modalities via a learned projection:
\begin{equation}
    \mathbf{e}_v = \text{MLP}([\mathbf{e}_v^{\text{text}} \| \mathbf{e}_v^{\text{struct}} \| \mathbf{e}_v^{\text{visual}}])
\end{equation}
where $\|$ denotes concatenation. The MLP is trained end-to-end with the retrieval objective.

\subsection{Temporal DOM Knowledge Graph}
\label{sec:temporal}

Given a sequence of DOM snapshots $\{T_1, T_2, \ldots, T_n\}$ of the same URL at times $\{t_1, t_2, \ldots, t_n\}$, we construct a temporal knowledge graph $\mathcal{G}^{\text{temp}}$.

\paragraph{Step 1: DOM Alignment.}
For consecutive snapshots $(T_i, T_{i+1})$, we compute the tree edit distance using the Zhang-Shasha algorithm~\cite{zhangshasha1989}, which produces an optimal mapping between nodes in the two trees. This mapping identifies:
\begin{itemize}
    \item \textbf{Unchanged nodes}: Same content and position
    \item \textbf{Modified nodes}: Same position, different content (e.g., price changed)
    \item \textbf{Inserted nodes}: New nodes in $T_{i+1}$
    \item \textbf{Deleted nodes}: Nodes removed from $T_i$
\end{itemize}

\paragraph{Step 2: Temporal Triple Extraction.}
For each modified node, we extract a temporal knowledge triple:
\begin{equation}
    (e, r, v_{\text{old}}, t_i) \rightarrow (e, r, v_{\text{new}}, t_{i+1})
\end{equation}
where $e$ is the entity (identified by DOM ancestry path), $r$ is the relation (inferred from tag semantics and CSS class), $v_{\text{old}}$ and $v_{\text{new}}$ are the old and new values, and $t_i, t_{i+1}$ are timestamps.

For example, if a \texttt{<span class="price">} node changes from ``\$99.99'' to ``\$79.99'':
\begin{equation}
    (\text{Product\_X}, \text{price}, \$99.99, t_1) \rightarrow (\text{Product\_X}, \text{price}, \$79.99, t_2)
\end{equation}

\paragraph{Step 3: Temporal Graph Construction.}
We organize temporal triples into a bi-level graph inspired by TG-RAG~\cite{tgrag2025}:
\begin{itemize}
    \item \textbf{Entity-level graph}: Nodes are entities with edges representing relationships (e.g., product $\rightarrow$ price, product $\rightarrow$ availability)
    \item \textbf{Time-level graph}: A hierarchical time index (day $\rightarrow$ week $\rightarrow$ month) with temporal summaries at each level
\end{itemize}

\subsection{Confidence-Guided Tree Traversal}
\label{sec:retrieval}

Given a user query $q$, we retrieve relevant DOM context via top-down tree traversal rather than flat top-$k$ retrieval.

\begin{algorithm}[t]
\caption{Confidence-Guided Tree Traversal}
\label{alg:traversal}
\begin{algorithmic}[1]
\REQUIRE Query embedding $\mathbf{q}$, DOM-KG root $r$, confidence threshold $\tau_h$, $\tau_l$, max tokens $L$
\STATE $\text{context} \leftarrow \emptyset$, $\text{queue} \leftarrow \{r\}$
\WHILE{$\text{queue} \neq \emptyset$ \AND $|\text{context}| < L$}
    \STATE $v \leftarrow \text{queue.pop()}$
    \STATE $s_v \leftarrow \text{sim}(\mathbf{q}, \mathbf{e}_v)$ \COMMENT{Relevance score}
    \IF{$s_v > \tau_h$ \AND $v$ has children}
        \STATE \COMMENT{High confidence: drill deeper}
        \FOR{each child $c$ of $v$}
            \STATE $\text{queue.push}(c)$
        \ENDFOR
    \ELSIF{$s_v > \tau_l$}
        \STATE \COMMENT{Medium confidence: retrieve whole subtree}
        \STATE $\text{context} \leftarrow \text{context} \cup \text{html}(v)$
    \ELSE
        \STATE \COMMENT{Low confidence: prune}
    \ENDIF
\ENDWHILE
\RETURN $\text{context}$
\end{algorithmic}
\end{algorithm}

The thresholds $\tau_h$ and $\tau_l$ are learned on a validation set. This approach achieves $O(b \cdot \log n)$ complexity where $b$ is the average branching factor and $n$ is the number of nodes, compared to $O(n)$ for flat retrieval.

\subsection{Cross-Site DOM Fingerprinting}
\label{sec:fingerprint}

To enable zero-shot extraction on unseen websites, we learn DOM subtree fingerprints via contrastive learning.

\paragraph{Subtree Representation.}
For a DOM subtree rooted at node $v$, we compute a fingerprint by aggregating tri-modal embeddings of all descendant nodes using a Graph Attention Network (GAT):
\begin{equation}
    \mathbf{f}_v = \text{GAT}(\{\mathbf{e}_u : u \in \text{subtree}(v)\}, \mathcal{E}_{\text{subtree}(v)})
\end{equation}

\paragraph{Contrastive Objective.}
We train with InfoNCE loss where:
\begin{itemize}
    \item \textbf{Positive pairs}: Subtrees containing the same semantic element (e.g., ``product price'') from different websites
    \item \textbf{Negative pairs}: Subtrees containing different semantic elements (e.g., ``price'' vs. ``navigation menu'')
\end{itemize}
\begin{equation}
    \mathcal{L} = -\log \frac{\exp(\text{sim}(\mathbf{f}_i, \mathbf{f}_j^+) / \tau)}{\sum_{k} \exp(\text{sim}(\mathbf{f}_i, \mathbf{f}_k) / \tau)}
\end{equation}

At inference time, given a new website, we compute fingerprints for all subtrees and match them against learned prototypes to identify semantic elements (price, name, availability, etc.) without any site-specific training.

% ============================================================
\section{Experimental Setup}
\label{sec:experiments}

\subsection{Datasets}

We evaluate on three categories of tasks:

\paragraph{Static HTML QA.} Following HtmlRAG~\cite{htmlrag2025}, we use six QA datasets: ASQA, HotpotQA, NQ, TriviaQA, MuSiQue, and ELI5, with HTML documents retrieved via Bing search.

\paragraph{Web Extraction.} We use the SWDE dataset~\cite{swde} (124K pages, 8 verticals) and evaluate zero-shot cross-site transfer.

\paragraph{Temporal Web QA.} We introduce \textbf{TempWebQA}, a new benchmark of product pages crawled at multiple time points from Common Crawl snapshots, with annotated temporal questions about price changes, availability, and sales events.

\subsection{Baselines}

\todo{List all baselines with citations}

\subsection{Implementation Details}

\todo{Model sizes, hyperparameters, training details}

% ============================================================
\section{Results}
\label{sec:results}

\todo{Main results tables}

\subsection{Static HTML QA}
\todo{Table comparing with HtmlRAG and baselines on 6 datasets}

\subsection{Web Extraction}
\todo{Table comparing with AXE, SCRIBES on SWDE}

\subsection{Temporal Web QA}
\todo{Table on TempWebQA benchmark}

\subsection{Retrieval Efficiency}
\todo{Latency and token count comparison}

\subsection{Ablation Studies}
\todo{Modality ablation, temporal ablation, retrieval strategy ablation}

% ============================================================
\section{Analysis}
\label{sec:analysis}

\subsection{Impact of Visual Grounding}
\todo{Case studies showing where visual position matters}

\subsection{Temporal Reasoning Examples}
\todo{Examples of temporal QA with provenance chains}

\subsection{Cross-Site Transfer}
\todo{Analysis of fingerprint quality across site families}

\subsection{Error Analysis}
\todo{Failure cases and limitations}

% ============================================================
\section{Conclusion}
\label{sec:conclusion}

We presented \ours{}, a novel RAG framework that treats HTML DOM trees as temporally-evolving, visually-grounded knowledge graphs. By introducing tri-modal DOM node embeddings, temporal knowledge graph construction via structural tree differencing, confidence-guided tree traversal for retrieval, and cross-site DOM fingerprinting, our system addresses fundamental limitations of existing HTML-based RAG approaches. Experiments demonstrate improvements across static QA, web extraction, and temporal QA tasks. We release our code, models, and the TempWebQA benchmark to facilitate future research.

\paragraph{Future Work.} We envision extending \ours{} to support multi-page agentic reasoning (following links across pages), integration with web agents for interactive information gathering, and application to domains beyond e-commerce such as news monitoring and regulatory compliance.

% ============================================================
% References
\begin{thebibliography}{99}

\bibitem{htmlrag2025}
J.~Tan, Z.~Dou, W.~Wang, M.~Wang, W.~Chen, and J.-R.~Wen.
\newblock {HtmlRAG}: {HTML} is Better Than Plain Text for Modeling Retrieved Knowledge in {RAG} Systems.
\newblock In \emph{Proceedings of the ACM Web Conference (WWW)}, 2025.

\bibitem{markuplm2022}
J.~Li, Y.~Xu, L.~Cui, and F.~Wei.
\newblock {MarkupLM}: Pre-training of Text and Markup Language for Visually-rich Document Understanding.
\newblock In \emph{Proceedings of ACL}, 2022.

\bibitem{domlm2022}
X.~Deng, P.~Shiralkar, C.~Lockard, B.~Huang, and H.~Sun.
\newblock {DOM-LM}: Learning Generalizable Representations for {HTML} Documents.
\newblock \emph{arXiv preprint arXiv:2201.10608}, 2022.

\bibitem{webformer2022}
Q.~Wang, Y.~Fang, A.~Ravula, F.~Feng, X.~Quan, and D.~Liu.
\newblock {WebFormer}: The Web-page Transformer for Structure Information Extraction.
\newblock In \emph{Proceedings of the ACM Web Conference (WWW)}, 2022.

\bibitem{klarna2024}
A.~Hotti et al.
\newblock The {Klarna} Product Page Dataset: Web Element Nomination with Graph Neural Networks and Large Language Models.
\newblock In \emph{OpenReview}, 2024.

\bibitem{axe2025}
Low-Cost Cross-Domain Web Structured Information Extraction.
\newblock \emph{arXiv preprint arXiv:2602.01838}, 2025.

\bibitem{scribes2025}
Web-Scale Script-Based Semi-Structured Data Extraction with Reinforcement Learning.
\newblock \emph{arXiv preprint arXiv:2510.01832}, 2025.

\bibitem{prune4web2025}
{DOM} Tree Pruning Programming for Web Agent.
\newblock \emph{arXiv preprint arXiv:2511.21398}, 2025.

\bibitem{pinterest2025}
Cross-Domain Web Information Extraction at Pinterest.
\newblock \emph{arXiv preprint arXiv:2508.01096}, 2025.

\bibitem{colpali2024}
{ColPali}: Efficient Document Retrieval with Vision Language Models.
\newblock \emph{arXiv preprint arXiv:2407.01449}, 2024.

\bibitem{snappy2024}
Spatially-Grounded Document Retrieval via Patch-to-Region Relevance Propagation.
\newblock \emph{arXiv preprint arXiv:2512.02660}, 2024.

\bibitem{scan2025}
{SCAN}: Semantic Document Layout Analysis for Textual and Visual {RAG}.
\newblock \emph{arXiv preprint arXiv:2505.14381}, 2025.

\bibitem{tgrag2025}
Time-Sensitive Modeling and Retrieval for Evolving Knowledge.
\newblock \emph{arXiv preprint arXiv:2510.13590}, 2025.

\bibitem{starrag2025}
Right Answer at the Right Time --- Temporal {RAG} via Graph Summarization.
\newblock \emph{arXiv preprint arXiv:2510.16715}, 2025.

\bibitem{streamingrag2025}
{StreamingRAG}: Real-time Contextual Retrieval and Generation Framework.
\newblock \emph{arXiv preprint arXiv:2501.14101}, 2025.

\bibitem{evokg2025}
Temporal Reasoning over Evolving Knowledge Graphs.
\newblock \emph{arXiv preprint arXiv:2509.15464}, 2025.

\bibitem{atom2025}
{ATOM}: Adaptive and Optimized Dynamic Temporal Knowledge Graph Construction using {LLMs}.
\newblock \emph{arXiv preprint arXiv:2510.22590}, 2025.

\bibitem{hdna2023}
{HDNA}: A Graph-based Change Detection in {HTML} Pages.
\newblock \emph{arXiv preprint arXiv:2310.03891}, 2023.

\bibitem{zhangshasha1989}
K.~Zhang and D.~Shasha.
\newblock Simple Fast Algorithms for the Editing Distance between Trees and Related Problems.
\newblock \emph{SIAM Journal on Computing}, 18(6):1245--1262, 1989.

\bibitem{vdiff2004}
G.~Cobéna, S.~Abiteboul, and A.~Marian.
\newblock Detecting Changes in {XML} Documents.
\newblock In \emph{Proceedings of ICDE}, 2002.

\bibitem{langchain2022}
H.~Chase.
\newblock {LangChain}, 2022.

\bibitem{llamaindex2022}
J.~Liu.
\newblock {LlamaIndex}, 2022.

\bibitem{swde}
Q.~Hao, R.~Cai, Y.~Pang, and L.~Zhang.
\newblock From One Tree to a Forest: a Unified Solution for Structured Web Data Extraction.
\newblock In \emph{Proceedings of SIGIR}, 2011.

\bibitem{rag_survey2025}
A Comprehensive Survey of Retrieval-Augmented Generation ({RAG}).
\newblock \emph{arXiv preprint arXiv:2506.00054}, 2025.

\bibitem{productkg2025}
{AI} Agent-Driven Framework for Automated Product Knowledge Graph Construction in E-Commerce.
\newblock \emph{arXiv preprint arXiv:2511.11017}, 2025.

\bibitem{webappkg2024}
Representing Web Applications As Knowledge Graphs.
\newblock \emph{arXiv preprint arXiv:2410.17258}, 2024.

\end{thebibliography}

\end{document}
